\documentclass[11pt]{article}
\usepackage[margin=1in]{geometry}
\usepackage{graphicx}
\usepackage{booktabs}
\usepackage{longtable}
\usepackage{hyperref}
\usepackage{listings}
\usepackage{xcolor}
\lstset{basicstyle=\ttfamily\footnotesize, breaklines=true, columns=fullflexible}
\title{Independent Replication of Khoder et al. 2022:\\ Whole-Genome Identification of \emph{Neisseria} spp.\\ (BVBRC-121)}
\author{OpenClaw agent (bvbrc-121 subagent, run 2026-07-05)}
\date{\today}
\begin{document}
\maketitle

\begin{abstract}
Independent replication of Khoder \emph{et al.} 2022 (PMID 36362240, DOI 10.3390/ijms232113456), which used OrthoANI, isDDH, and Roary pangenome analysis to reclassify four Lebanese Neisseria clinical isolates that had been misidentified as \emph{N.\ gonorrhoeae} by the API-NH biochemical panel. Using an all-free-endpoint workflow (NCBI \texttt{datasets} $\to$ \texttt{skani} ANI $\to$ UPGMA tree $\to$ Argo \texttt{gpt-5.2} LLM-judge) on 19 genomes (4 Lebanese isolates + 15 stratified references), we reproduce the paper's principal diagnostic message (all four isolates $\leq$85.2\% ANI to \emph{N.\ gonorrhoeae} FA1090, well below the 95\% species cutoff), independently confirm R20 as \emph{N.\ mucosa} (96.06\% ANI), and independently corroborate the paper's collateral claim that NCBI contains taxonomic errors by finding that the deposited ``\emph{N.\ sicca} VK64'' assembly (GCF\_000260655.1) is 96.83\% ANI to \emph{N.\ macacae}. However, the specific \emph{N.\ flavescens} assignment for R19, R21, and R23 does \emph{not} survive a strict 95\% ANI test against modern type-strain references: those isolates are actually closer to \emph{N.\ subflava} (98.04\%) and \emph{N.\ perflava} (98.05\%) than to any labeled \emph{N.\ flavescens} reference. This reinforces the paper's own broader claim of taxonomic ambiguity in the flavescens/perflava/subflava/mucosa complex. \textbf{Overall verdict: PARTIAL} (coverage 60\%, agreement 75\%, LLM-judge confidence high).
\end{abstract}

\section{Paper summary}

Khoder \emph{et al.} sequenced four Neisseria isolates (R19, R20, R21, R23) from semen samples of infertile Lebanese men. API-NH called all four \emph{N.\ gonorrhoeae}; MALDI-TOF reassigned R19/R21/R23 to \emph{N.\ flavescens} and R20 to \emph{N.\ mucosa}; 16S rDNA resolved only to genus. The authors then draft-sequenced (Illumina MiSeq, A5 pipeline, Prokka + RAST) and compared against 128 NCBI reference genomes using OrthoANI (ezbiocloud), isDDH (GGDC formula~2), and Roary pangenome (Galaxy Australia). Cutoffs: 95\% ANI / 70\% dDDH.

\section{Claims}

\begin{longtable}{@{}p{0.4cm}p{7.5cm}p{2.2cm}p{2.5cm}p{2.5cm}@{}}
\toprule
ID & Claim & Type & Tested here? & Result \\ \midrule \endhead
C1 & R19/R20/R21/R23 are NOT \emph{N.\ gonorrhoeae} & reclassification & Yes & \textbf{Reproduced} \\
C2 & R19/R21/R23 = \emph{N.\ flavescens}; R20 = \emph{N.\ mucosa} & species assign. & Yes & \textbf{Partial} (R20 yes; R19/R21/R23 fail strict 95\% cutoff) \\
C3 & isDDH + OrthoANI + pangenome combined is best & methodology & No (ANI only) & Not tested \\
C4 & NCBI Neisseria has many taxonomic errors & database quality & Yes & \textbf{Reproduced} (VK64 is macacae, not sicca) \\
C5 & Genuine ambiguity in flavescens/perflava/subflava complex & biology & Yes & \textbf{Reproduced with new evidence} \\
\bottomrule
\end{longtable}

\section{Method}

\begin{enumerate}
\item Fetched paper PDF from Europe PMC (\texttt{europepmc.org/articles/PMC9657967?pdf=render}); extracted text with local \texttt{pdftotext} (poppler). Parsed Methods and Table 1 for the four Lebanese GenBank accessions.
\item Downloaded 4 Lebanese isolates + 15 stratified references with \texttt{datasets download genome accession} (NCBI datasets CLI 18.32.0) on uicgpu via HTTP proxy. Every FASTA passed a header-Neisseria sanity check (caught 7 wrong-taxon initial guesses --- fixed by taxon-search).
\item All-vs-all ANI with \texttt{skani dist --min-af 0 -s 70}. skani is the state-of-the-art fastANI/OrthoANI replacement (Shaw \& Yu, \emph{Nature Methods} 2023); correlates $>0.99$ with OrthoANI at 90--100\% ANI.
\item Cross-check with \texttt{mash sketch -k 21 -s 10000} + \texttt{mash dist} --- qualitatively consistent.
\item UPGMA tree from ANI distances via \texttt{scipy.cluster.hierarchy.linkage(method='average')}; species partition at cutoff 5.0 (equivalent to 95\% ANI); Newick export by manual traversal of the scipy linkage matrix; heatmap+dendrogram with matplotlib.
\item Per-claim verification script produces \texttt{claim\_verification.json}; then LLM-judge (Argo \texttt{gpt-5.2} via cherryrd LiteLLM aggregator \texttt{<tailnet-aggregator>:4000/v1}, free endpoint; \texttt{claude-opus-4.8} was 502-ing) evaluates coverage, agreement, per-claim tested/reproduced booleans, and issues verdict.
\end{enumerate}

\section{Results vs paper}

Full ANI matrix in \texttt{report/evidence/ani\_matrix\_final.tsv}. Excerpt (rows = Lebanese isolates + VK64 misclass check; cols = key references; boldface crosses 95\% species boundary):

\begin{longtable}{@{}lrrrrrrrrr@{}}
\toprule
 & R19 & R20 & R21 & R23 & FA1090 & ATCC13120 & ATCC49275 & C2008000159 & ATCC33926 \\
 & & & & & (gono) & (flav) & (subf) & (mucosa) & (macacae) \\ \midrule
R19 & 100.00 & 84.41 & \textbf{95.50} & 94.32 & 84.57 & 94.08 & 94.35 & 84.56 & 84.18 \\
R20 & 84.41 & 100.00 & 83.86 & 85.22 & 85.21 & 85.39 & 86.35 & \textbf{96.06} & 93.64 \\
R21 & \textbf{95.50} & 83.86 & 100.00 & 94.58 & 83.63 & 93.98 & 94.62 & 84.04 & 84.45 \\
R23 & 94.32 & 85.22 & 94.58 & 100.00 & 84.54 & 94.56 & \textbf{96.50} & 86.06 & 85.98 \\
VK64 & 84.22 & 93.16 & 84.08 & 85.19 & 85.83 & 85.29 & 86.02 & 93.66 & \textbf{96.51--96.83} \\
\bottomrule
\end{longtable}

\subsection{Per-claim outcome}

\textbf{C1 --- Reproduced.} All four Lebanese isolates: 83.63--85.22\% ANI vs \emph{N.\ gonorrhoeae} FA1090 (Cluster A) and 84.15--85.48\% vs \emph{N.\ meningitidis} MC58 (Cluster A). API-NH mis-call decisively contradicted.

\textbf{C2 --- Partial.} R20 vs \emph{N.\ mucosa} C2008000159 = 96.06\% $\Rightarrow$ mucosa confirmed. R19/R21/R23 vs the closest labeled \emph{N.\ flavescens} reference (ATCC13120 type-strain equivalent) are only 93.98--94.56\%, \emph{below} the 95\% species cutoff. R19 is 98.04\% to \emph{N.\ subflava} C2005001510 and 98.05\% to \emph{N.\ perflava} 27098\_8\_142 --- closer to those species than to any flavescens ref. R23 is 96.50\% to \emph{N.\ subflava} ATCC49275 (type strain, complete chromosome). R19/R23 and R21/R23 pairwise ANIs sit at 94.32/94.58\% --- below the paper's own cutoff for a same-species assignment. The paper's flavescens label is \emph{unstable} under a strict $\geq$95\% ANI closest-species rule.

\textbf{C3 --- Not tested.} GGDC isDDH is web-only (no batch API); Roary pangenome would require Prokka annotation of all 19 genomes plus a Roary run --- out of subagent time budget.

\textbf{C4 --- Reproduced.} GCF\_000260655.1 is deposited in NCBI as ``\emph{N.\ sicca} VK64''. Its top ANI hits are 96.83\% and 96.51\% to \emph{N.\ macacae} ATCC33926 (two different assemblies of the same strain), both above 95\%. VK64 should be reclassified as \emph{N.\ macacae}. This is a direct independent example of the mislabeling problem the paper identifies.

\textbf{C5 --- Reproduced with new evidence.} Our matrix produces a $\geq$95\% ANI cluster spanning R19, R21, CD-NF2 (labeled flavescens), UMB0210 / 27098 (labeled perflava), and C2005001510 (labeled subflava) --- a five-way, three-species mixed cluster that violates the notion of species as a discrete unit under the 95\% ANI rule. This dramatically supports the paper's argument that stronger cutoffs are needed.

\subsection{Tree topology}

See \texttt{report/evidence/tree\_heatmap.png}. Major clades at the 95\% ANI cutoff:

\begin{itemize}
\item \textbf{Cluster A} (flavescens/perflava/subflava commensal): R19, R21, R23, plus 3 flavescens refs, 2 perflava refs, 2 subflava refs.
\item \textbf{Cluster B} (mucosa/macacae commensal): R20, mucosa C2008000159, macacae ATCC33926 (both assemblies), the mislabeled ``\emph{N.\ sicca}'' VK64.
\item \textbf{Cluster C} (pathogenic): \emph{N.\ gonorrhoeae}, \emph{N.\ meningitidis}, \emph{N.\ lactamica}.
\item \textbf{Outgroup}: \emph{N.\ elongata} NCTC10660.
\end{itemize}

Cluster-level topology agrees with the paper's Figures 1--2.

\section{Verdict + justification}

\textbf{Verdict: PARTIAL} (coverage 60\%, agreement 75\%, LLM-judge confidence high; judge model \texttt{argo:gpt-5.2}).

The paper's diagnostic headline reproduces cleanly, the mucosa-for-R20 assignment reproduces, and the ``NCBI has mislabels'' collateral claim is directly independently confirmed. However, the \emph{N.\ flavescens} label for R19/R21/R23 does not survive a strict 95\% ANI test, and this is not a technical failure of the replication --- it is a fragility of the paper's own conclusion under its own stated criteria. The paper's broader C5 claim about taxonomic ambiguity in the flavescens/perflava/subflava complex is strongly reinforced. Two of the paper's three methodological pillars (isDDH via GGDC web tool, Roary pangenome on Galaxy) were not rerun; ANI alone is the primary numerical evidence here.

\section{Open questions}

Five NEW open questions arose from this rerun; full detail (with \texttt{basis} and \texttt{next\_steps}) in \texttt{report/open\_questions.json}.

\begin{enumerate}
\item[\textbf{Q1}] Does R19/R21/R23 straddling the 95\% ANI cutoff persist at the paper's full 128-reference scale, or does denser sampling collapse them onto \emph{N.\ flavescens}?
\item[\textbf{Q2}] What is the base-rate of $>$95\%-ANI-mismatch-to-deposited-label across all NCBI RefSeq Neisseria? A systematic audit would turn the paper's anecdote into a quantified curation problem.
\item[\textbf{Q3}] Is the low R19/R23 pairwise 94.32\% ANI a real biological signal or a draft-assembly artifact (contig counts 34 vs 79)? A reassembly from raw reads would resolve this.
\item[\textbf{Q4}] The paper never states the decision rule when ``closest species'' and ``strict 95\%'' disagree. Which rule was used, and what would consistent application do to the paper's assignments?
\item[\textbf{Q5}] How often do isDDH, ANI, and pangenome actually agree on the 19-genome test set (four-tool concordance study, Fleiss' kappa across tools)?
\end{enumerate}

\section*{Reproduction}

Everything is in \texttt{work/}: \texttt{fetch\_*.sh} (three-pass genome download, with lessons applied), \texttt{analyze\_ani.py} (matrix + tree + heatmap + per-claim JSON), \texttt{llm\_judge.py} (Argo aggregator, \texttt{gpt-5.2} model). Wall-clock for a clean re-run: $\sim$10 minutes. Zero GPU time.

\end{document}
